\begin{exercisechunk}
\item \textbf{Permanent state and finite-window prediction.}
\label[exercise]{ex:l5-memory}
Use the process in \cref{eq:constant-state}.
\begin{enumerate}[(a),beginpenalty=10000]
\item Complete the square in the posterior density of \(U\) given
\(Z_{1:n}\) to derive \cref{eq:constant-posterior}.
Give a recursion that stores only a running sum and a count.
\item Derive the minimum prediction risks using all \(n\) observations
and only the last \(p\le n\). Prove that the latter bound applies to
nonlinear predictors as well as linear AR predictors.
\item Derive the expected conditional KL in
\cref{prop:memory-advantage} directly from the Gaussian KL formula.
You may use conditional-expectation orthogonality.
\item Explain why recovering \(U\) accurately along one infinite
trajectory does not consistently estimate an unknown prior variance
\(\tau^2\). What additional data would supply repeated draws of \(U\)?
\end{enumerate}
\end{exercisechunk}
